\documentclass{article}%
\usepackage[T1]{fontenc}%
\usepackage[utf8]{inputenc}%
\usepackage{lmodern}%
\usepackage{textcomp}%
\usepackage{lastpage}%
\usepackage{geometry}%
\geometry{margin=1in,headheight=14pt,headsep=25pt}%
\usepackage{hyperref}%
\usepackage{enumitem}%
\usepackage{fancyhdr}%
\usepackage{xcolor}%
\usepackage{url}%
\usepackage{breakurl}%
%

            \hypersetup{
                colorlinks=true,
                linkcolor=blue,
                filecolor=magenta,
                urlcolor=blue
            }
            \pagestyle{fancy}
            \fancyhf{}
            \rhead{Generated on \today}
            \cfoot{\thepage}
            
            % Configure enumeration settings
            \setlist[enumerate,1]{label=\arabic*.}
            \setlist[enumerate,2]{label=\alph*.}
            \setlist[enumerate,3]{label=\Alph*.}
            \setlist[enumerate]{itemsep=0.5em}
            
            % Define a command for red text
            \newcommand{\incorrect}[1]{\textcolor{red}{#1}}
        %
\lhead{2024{-}10{-}22{-}Convexity}%
\title{Summary for 2024{-}10{-}22{-}Convexity}%
\author{Generated by Scribe.AI}%
\date{\today}%
%
\begin{document}%
\normalsize%
\maketitle%
\section{Summary}%
\label{sec:Summary}%
This lecture on convexity in linear programming introduced fundamental concepts and theorems related to convex sets and their properties.  The core idea revolves around the geometric properties of feasible regions in optimization problems.

* **Convex Set:** A set S in ℝ<sup>m</sup> is convex if, for any two points within the set, the line segment connecting them is entirely contained within the set.

* **Convex Combination:** A convex combination of points z<sub>1</sub>, z<sub>2</sub>,..., z<sub>n</sub> in ℝ<sup>m</sup> is a weighted sum  t<sub>1</sub>z<sub>1</sub> + t<sub>2</sub>z<sub>2</sub> + ... + t<sub>n</sub>z<sub>n</sub>, where each t<sub>i</sub> is non-negative and their sum equals 1.

* **Equivalence of Convexity and Convex Combinations:** A set is convex if and only if it contains all convex combinations of its points.  This theorem is proven using induction, demonstrating that if a set is convex, it must include all such combinations, and vice-versa.

* **Convex Hull:** The convex hull of a set S, denoted Conv(S), is the smallest convex set containing S, equivalent to the intersection of all convex sets containing S, and also equivalent to the set of all convex combinations of finitely many points from S.

* **Caratheodory's Theorem:** The convex hull of a set S in ℝ<sup>m</sup> consists of all convex combinations of at most m+1 points from S.  This is proven using linear programming, showing that an optimal solution to a specific linear program will always involve at most m+1 points.

* **Separation Theorem:** Two disjoint nonempty polyhedra in ℝ<sup>n</sup> can be separated by two disjoint half-spaces. This theorem is proven using Farkas' Lemma, which provides a condition for the non-existence of solutions to a system of inequalities.

* **Farkas' Lemma:** The system of inequalities Ax ≤ b has no solution if and only if there exists a vector y such that A<sup>T</sup>y = 0, y ≥ 0, and b<sup>T</sup>y < 0.  This lemma is fundamental to linear programming duality and is used in the proof of the Separation Theorem.

%
\end{document}